\documentclass[12pt]{article}
\usepackage[UTF8]{inputenc}
\usepackage{thaisupp}
\usepackage{enumitem}
\begin{document}
\title{ตัวเก็บประจุ}
\maketitle
\section*{In-class Questions}
\begin{enumerate}[label=\arabic*.]
\item ตัวเก็บประจุ (Capacitor) คืออุปกรณ์ที่ทำหน้าที่อะไร? \hfill \textbf{\small (Main)}
\begin{enumerate}[label=)]
\item ก) เก็บประจุไฟฟ้า
\item ข) เก็บพลังงานไฟฟ้าในรูปสนามแม่เหล็ก
\item ค) เปลี่ยนไฟฟ้ากระแสสลับเป็นไฟฟ้ากระแสตรง
\item ง) ขยายสัญญาณไฟฟ้า
\end{enumerate}
\item ความจุไฟฟ้า C = Q/V มีหน่วยเป็นอะไร? \hfill \textbf{\small (Main)}
\begin{enumerate}[label=)]
\item ก) โวลต์ต่อคูลอมบ์ (V/C)
\item ข) คูลอมบ์ต่อโวลต์ (C/V)
\item ค) จูลต่อคูลอมบ์ (J/C)
\item ง) แอมแปร์ต่อวินาที (A/s)
\end{enumerate}
\item ตัวเก็บประจุแผ่นขนานมีความจุ C = ε₀A/d ถ้าเพิ่มพื้นที่แผ่นเป็น 2 เท่า โดยระยะห่างคงเดิม ความจุจะเปลี่ยนเป็นเท่าใด? \hfill \textbf{\small (Main)}
\begin{enumerate}[label=)]
\item ก) C/2
\item ข) C
\item ค) 2C
\item ง) 4C
\end{enumerate}
\item พลังงานที่เก็บในตัวเก็บประจุมีสูตรใด? \hfill \textbf{\small (Main)}
\begin{enumerate}[label=)]
\item ก) U = QV
\item ข) U = ½QV
\item ค) U = ½CV²
\item ง) ถูกทั้งข้อ ข และ ค
\end{enumerate}
\item ตัวเก็บประจุสองตัว C₁ = 2 μF และ C₂ = 3 μF ต่ออนุกรมกัน ความจุเทียบเท่าเป็นเท่าใด? \hfill \textbf{\small (Main)}
\begin{enumerate}[label=)]
\item ก) 1 μF
\item ข) 1.2 μF
\item ค) 5 μF
\item ง) 6 μF
\end{enumerate}
\item ตัวเก็บประจุสองตัว C₁ = 2 μF และ C₂ = 3 μF ต่อขนานกัน ความจุเทียบเท่าเป็นเท่าใด? \hfill \textbf{\small (Main)}
\begin{enumerate}[label=)]
\item ก) 1 μF
\item ข) 1.2 μF
\item ค) 5 μF
\item ง) 6 μF
\end{enumerate}
\item ตัวเก็บประจุ 10 μF ต่อกับความต่างศักย์ 100 V พลังงานที่เก็บได้เป็นเท่าใด? \hfill \textbf{\small (Main)}
\begin{enumerate}[label=)]
\item ก) 0.05 J
\item ข) 0.5 J
\item ค) 5 J
\item ง) 50 J
\end{enumerate}
\item ถ้าตัวเก็บประจุ 100 μF ต่อกับความต่างศักย์ 50 V จะเก็บประจุได้กี่คูลอมบ์? \hfill \textbf{\small (Main)}
\begin{enumerate}[label=)]
\item ก) 2×10⁻³ C
\item ข) 5×10⁻³ C
\item ค) 2×10⁻² C
\item ง) 5×10⁻² C
\end{enumerate}
\item ตัวเก็บประจุ 5 μF มีประจุ 2×10⁻⁴ C พลังงานที่เก็บเป็นเท่าใด? \hfill \textbf{\small (Main)}
\begin{enumerate}[label=)]
\item ก) 2×10⁻³ J
\item ข) 4×10⁻³ J
\item ค) 8×10⁻³ J
\item ง) 1×10⁻² J
\end{enumerate}
\item ไดอิเล็กทริก (Dielectric) ในตัวเก็บประจุแผ่นขนานมีหน้าที่อะไร? \hfill \textbf{\small (Main)}
\begin{enumerate}[label=)]
\item ก) เพิ่มความต้านทาน
\item ข) เพิ่มความจุไฟฟ้า
\item ค) ลดความจุไฟฟ้า
\item ง) เปลี่ยนไฟฟ้ากระแสตรงเป็นสลับ
\end{enumerate}
\item ตัวเก็บประจุสามตัว 2 μF แต่ละตัว ต่ออนุกรมกัน ความจุเทียบเท่าเป็นเท่าใด? \hfill \textbf{\small (Main)}
\begin{enumerate}[label=)]
\item ก) 6 μF
\item ข) 3 μF
\item ค) 2 μF
\item ง) 2/3 μF
\end{enumerate}
\item ตัวเก็บประจุสามตัว 3 μF แต่ละตัว ต่อขนานกัน ความจุเทียบเท่าเป็นเท่าใด? \hfill \textbf{\small (Main)}
\begin{enumerate}[label=)]
\item ก) 1 μF
\item ข) 3 μF
\item ค) 6 μF
\item ง) 9 μF
\end{enumerate}
\item ตัวเก็บประจุแผ่นขนานมีพื้นที่ 0.05 m² ระยะห่าง 0.001 m มีไดอิเล็กทริก εᵣ = 4 ความจุเป็นเท่าใด? (ε₀ = 8.85×10⁻¹²) \hfill \textbf{\small (Main)}
\begin{enumerate}[label=)]
\item ก) 1.77×10⁻⁹ F
\item ข) 1.77×10⁻¹⁰ F
\item ค) 7.08×10⁻¹⁰ F
\item ง) 7.08×10⁻⁹ F
\end{enumerate}
\item ตัวเก็บประจุที่เก็บประจุ Q แล้วตัดออกจากวงจร ถ้าเพิ่มระยะห่างเป็น 2 เท่า พลังงานจะเปลี่ยนอย่างไร? \hfill \textbf{\small (Main)}
\begin{enumerate}[label=)]
\item ก) เพิ่มเป็น 2 เท่า
\item ข) เพิ่มเป็น 4 เท่า
\item ค) เท่าเดิม
\item ง) ลดลง 2 เท่า
\end{enumerate}
\item ตัวเก็บประจุ 10 μF ต่อกับความต่างศักย์ 200 V แล้วตัดออก แล้วต่อขนานกับตัวเก็บประจุ 30 μF ที่ไม่มีประจุ ความต่างศักย์ร่วมเป็นเท่าใด? \hfill \textbf{\small (Main)}
\begin{enumerate}[label=)]
\item ก) 25 V
\item ข) 50 V
\item ค) 100 V
\item ง) 200 V
\end{enumerate}
\item ตัวเก็บประจุแผ่นขนานเมื่อต่อกับแบตเตอรี่ ถ้าเพิ่มระยะห่างระหว่างแผ่นเป็น 2 เท่า โดยไม่ตัดแบตเตอรี่ ประจุจะเปลี่ยนอย่างไร? \hfill \textbf{\small (Main)}
\begin{enumerate}[label=)]
\item ก) เพิ่มเป็น 2 เท่า
\item ข) เท่าเดิม
\item ค) ลดลง 2 เท่า
\item ง) ลดลง 4 เท่า
\end{enumerate}
\item ตัวเก็บประจุ 4 μF ต่ออนุกรมกับกลุ่มตัวเก็บประจุ 6 μF และ 3 μF ที่ต่อขนานกัน ความจุเทียบเท่าเป็นเท่าใด? \hfill \textbf{\small (Main)}
\begin{enumerate}[label=)]
\item ก) 3 μF
\item ข) 5 μF
\item ค) 6 μF
\item ง) 13 μF
\end{enumerate}
\item ตัวเก็บประจุสองตัว 3 μF และ 6 μF ต่อขนานกับแบตเตอรี่ 12 V พลังงานรวมเป็นเท่าใด? \hfill \textbf{\small (Main)}
\begin{enumerate}[label=)]
\item ก) 4.32×10⁻⁴ J
\item ข) 8.64×10⁻⁴ J
\item ค) 1.296×10⁻³ J
\item ง) 2.16×10⁻³ J
\end{enumerate}
\item ตัวเก็บประจุแผ่นขนานต่อกับแบตเตอรี่ V เมื่อเปลี่ยนไดอิเล็กทริกจากอากาศ (εᵣ=1) เป็นแก้ว (εᵣ=8) ประจุเปลี่ยนจาก Q₀ เป็นเท่าใด? \hfill \textbf{\small (Main)}
\begin{enumerate}[label=)]
\item ก) Q₀
\item ข) 4Q₀
\item ค) 8Q₀
\item ง) Q₀/8
\end{enumerate}
\item ตัวเก็บประจุ C ต่อกับแบตเตอรี่ V เมื่อเพิ่มความต่างศักย์เป็น 2V พลังงานเพิ่มขึ้นเป็นกี่เท่า? \hfill \textbf{\small (Main)}
\begin{enumerate}[label=)]
\item ก) 2 เท่า
\item ข) 3 เท่า
\item ค) 4 เท่า
\item ง) 8 เท่า
\end{enumerate}
\item ตัวเก็บประจุแผ่นขนานมีความจุ C₀ เมื่อเปลี่ยนระยะห่างจาก d เป็น d/2 โดยประจุ Q คงที่ พลังงานเปลี่ยนเป็นเท่าใด? \hfill \textbf{\small (Main)}
\begin{enumerate}[label=)]
\item ก) U₀/2
\item ข) U₀
\item ค) 2U₀
\item ง) 4U₀
\end{enumerate}
\item ตัวเก็บประจุแผ่นขนานมีความจุ C ถ้าเพิ่มพื้นที่เป็น 3 เท่า และลดระยะห่างเป็น 1/2 เท่า ความจุจะเปลี่ยนเป็นเท่าใด? \hfill \textbf{\small (Main)}
\begin{enumerate}[label=)]
\item ก) 3C
\item ข) 6C
\item ค) 9C
\item ง) 12C
\end{enumerate}
\item ตัวเก็บประจุ C₁ มีประจุ Q ต่อขนานกับตัวเก็บประจุ C₂ ที่มีประจุ Q เช่นกัน พลังงานรวมก่อนและหลังต่อเป็นเท่าใด? \hfill \textbf{\small (Main)}
\begin{enumerate}[label=)]
\item ก) เท่าเดิม
\item ข) เพิ่มขึ้น
\item ค) ลดลง
\item ง) ขึ้นอยู่กับค่า C₁ และ C₂
\end{enumerate}
\item ตัวเก็บประจุมีความจุ 100 μF ระยะห่าง 1 mm ถ้าใส่แผ่นไดอิเล็กทริกหนา 0.5 mm เต็มความหนา ความจุจะเปลี่ยนอย่างไร? (εᵣ = 4) \hfill \textbf{\small (Main)}
\begin{enumerate}[label=)]
\item ก) เพิ่มเป็น 2 เท่า
\item ข) เพิ่มเป็น 4 เท่า
\item ค) เพิ่มเป็น 8 เท่า
\item ง) ไม่เปลี่ยน
\end{enumerate}
\item ตัวเก็บประจุสองตัว C เท่ากัน ต่ออนุกรมกัน แล้วต่อกับแบตเตอรี่ V ประจุบนตัวเก็บประจุแต่ละตัวเป็นเท่าใด? \hfill \textbf{\small (Main)}
\begin{enumerate}[label=)]
\item ก) CV/2
\item ข) CV
\item ค) 2CV
\item ง) V/(2C)
\end{enumerate}
\item ตัวเก็บประจุในวงจรไฟฟ้ากระแสตรงเมื่อต่อนานๆ จะทำหน้าที่อย่างไร? \hfill \textbf{\small (Extra)}
\begin{enumerate}[label=)]
\item ก) ยอมให้กระแสไหลผ่านได้สมบูรณ์
\item ข) ปิดกั้นกระแสไฟฟ้ากระแสตรง (หลังประจุเต็ม)
\item ค) ขยายสัญญาณ
\item ง) เปลี่ยนไฟฟ้ากระแสตรงเป็นสลับ
\end{enumerate}
\item 1 Farad มีค่าเท่ากับกี่ไมโครฟารัด (μF)? \hfill \textbf{\small (Extra)}
\begin{enumerate}[label=)]
\item ก) 10⁻⁶ μF
\item ข) 10⁶ μF
\item ค) 10³ μF
\item ง) 10⁻³ μF
\end{enumerate}
\item พลังงานในตัวเก็บประจุถูกเก็บในรูปใด? \hfill \textbf{\small (Extra)}
\begin{enumerate}[label=)]
\item ก) รูปกระแสไฟฟ้า
\item ข) รูปสนามแม่เหล็ก
\item ค) รูปสนามไฟฟ้าระหว่างแผ่น
\item ง) รูปความร้อน
\end{enumerate}
\item ตัวเก็บประจุแผ่นขนานเมื่อต่อกับแบตเตอรี่ ประจุบนแผ่นทั้งสองมีขนาดอย่างไร? \hfill \textbf{\small (Extra)}
\begin{enumerate}[label=)]
\item ก) เท่ากันแต่ชนิดเดียวกัน
\item ข) เท่ากันแต่ชนิดตรงข้ามกัน
\item ค) ขนาดต่างกันขึ้นอยู่กับพื้นที่
\item ง) ขนาดต่างกันขึ้นอยู่กับระยะห่าง
\end{enumerate}
\item ตัวเก็บประจุสองตัว 2 μF และ 4 μF ต่ออนุกรมกัน แล้วต่อกับแบตเตอรี่ 12 V ความต่างศักย์คร่อมตัวเก็บประจุ 2 μF เป็นเท่าใด? \hfill \textbf{\small (Extra)}
\begin{enumerate}[label=)]
\item ก) 2 V
\item ข) 4 V
\item ค) 8 V
\item ง) 12 V
\end{enumerate}
\item ตัวเก็บประจุสองตัว 3 μF และ 6 μF ต่ออนุกรมกัน ความจุรวมเป็นเท่าใด? \hfill \textbf{\small (Extra)}
\begin{enumerate}[label=)]
\item ก) 9 μF
\item ข) 3 μF
\item ค) 2 μF
\item ง) 1 μF
\end{enumerate}
\item ตัวเก็บประจุ 20 μF ต่อกับแบตเตอรี่ 50 V พลังงานที่เก็บได้เป็นเท่าใด? \hfill \textbf{\small (Extra)}
\begin{enumerate}[label=)]
\item ก) 2.5×10⁻² J
\item ข) 5.0×10⁻² J
\item ค) 1.0×10⁻¹ J
\item ง) 2.5×10⁻³ J
\end{enumerate}
\item ตัวเก็บประจุแผ่นขนานมีพื้นที่ A = 100 cm² ระยะห่าง d = 1 mm ความจุเป็นเท่าใด? (ε₀ = 8.85×10⁻¹² F/m) \hfill \textbf{\small (Extra)}
\begin{enumerate}[label=)]
\item ก) 8.85×10⁻¹¹ F
\item ข) 8.85×10⁻¹² F
\item ค) 8.85×10⁻¹³ F
\item ง) 8.85×10⁻¹⁰ F
\end{enumerate}
\item ตัวเก็บประจุ 10 μF ที่มีประจุ 10⁻³ C ถูกต่อขนานกับตัวเก็บประจุ 30 μF ที่ไม่มีประจุ พลังงานที่สูญเสียไปจากระบบเป็นเท่าใด? \hfill \textbf{\small (Extra)}
\begin{enumerate}[label=)]
\item ก) 1.5×10⁻³ J
\item ข) 3.75×10⁻³ J
\item ค) 5×10⁻³ J
\item ง) 7.5×10⁻³ J
\end{enumerate}
\item ตัวเก็บประจุแผ่นขนานเมื่อตัดออกจากวงจร (Q คงที่) แล้วเพิ่มระยะห่าง ความต่างศักย์จะเปลี่ยนอย่างไร? \hfill \textbf{\small (Extra)}
\begin{enumerate}[label=)]
\item ก) คงที่
\item ข) เพิ่มขึ้น
\item ค) ลดลง
\item ง) เป็นศูนย์
\end{enumerate}
\item ตัวเก็บประจุ 4 ตัว เท่ากัน C ต่อขนานกัน ความจุรวมเป็นเท่าใด? \hfill \textbf{\small (Extra)}
\begin{enumerate}[label=)]
\item ก) C/4
\item ข) C
\item ค) 2C
\item ง) 4C
\end{enumerate}
\item ตัวเก็บประจุ 2 μF และ 6 μF ต่ออนุกรมกัน แล้วต่อกับแบตเตอรี่ 24 V ประจุบนตัวเก็บประจุแต่ละตัวเป็นเท่าใด? \hfill \textbf{\small (Extra)}
\begin{enumerate}[label=)]
\item ก) 6 μC
\item ข) 18 μC
\item ค) 24 μC
\item ง) 36 μC
\end{enumerate}
\item ตัวเก็บประจุแผ่นขนานมีไดอิเล็กทริกพลอยกลาง (εᵣ = 6) แทรกเข้าไปครึ่งหนึ่งของช่องว่าง ความจุจะเปลี่ยนอย่างไร? (เทียบกับไม่มีไดอิเล็กทริก) \hfill \textbf{\small (Extra)}
\begin{enumerate}[label=)]
\item ก) เพิ่ม 3 เท่า
\item ข) เพิ่ม 4 เท่า
\item ค) เพิ่ม 6 เท่า
\item ง) เพิ่ม 1.5 เท่า
\end{enumerate}
\item ตัวเก็บประจุ 5 μF ต่อกับแบตเตอรี่ 100 V แล้วตัดออก จากนั้นต่อขนานกับตัวเก็บประจุ 15 μF ที่ไม่มีประจุ พลังงานหลังต่อเป็นเท่าใด? \hfill \textbf{\small (Extra)}
\begin{enumerate}[label=)]
\item ก) 1.25×10⁻² J
\item ข) 2.5×10⁻² J
\item ค) 3.75×10⁻² J
\item ง) 5×10⁻² J
\end{enumerate}
\item ตัวเก็บประจุชนิดใดที่มีความจุสูงและขนาดเล็ก เหมาะสำหรับใช้ในวงจรอิเล็กทรอนิกส์? \hfill \textbf{\small (Extra)}
\begin{enumerate}[label=)]
\item ก) ตัวเก็บประจุกระดาษ
\item ข) ตัวเก็บประจุแผ่นขนานอากาศ
\item ค) ตัวเก็บประจุเซรามิก/อิเล็กโทรไลติก
\item ง) ตัวเก็บประจุแก้ว
\end{enumerate}
\item ตัวเก็บประจุ 3 ตัว 2 μF, 3 μF, 6 μF ต่ออนุกรมกัน ความจุรวมเป็นเท่าใด? \hfill \textbf{\small (Extra)}
\begin{enumerate}[label=)]
\item ก) 1 μF
\item ข) 3 μF
\item ค) 6/11 μF
\item ง) 11/6 μF
\end{enumerate}
\item ตัวเก็บประจุ C ต่อกับแบตเตอรี่ V มีพลังงาน U ถ้าตัดแบตเตอรี่แล้วเพิ่มความจุเป็น 2C โดยไม่ให้ประจุรั่วไหล พลังงานจะเปลี่ยนเป็นเท่าใด? \hfill \textbf{\small (Extra)}
\begin{enumerate}[label=)]
\item ก) U/2
\item ข) U
\item ค) 2U
\item ง) 4U
\end{enumerate}
\item ตัวเก็บประจุแผ่นขนานต่อกับแบตเตอรี่ V ถ้าเพิ่มระยะห่างเป็น 3 เท่า โดยไม่ตัดแบตเตอรี่ พลังงานจะเปลี่ยนอย่างไร? \hfill \textbf{\small (Extra)}
\begin{enumerate}[label=)]
\item ก) เพิ่มเป็น 3 เท่า
\item ข) เพิ่มเป็น 9 เท่า
\item ค) ลดลง 3 เท่า
\item ง) ลดลง 9 เท่า
\end{enumerate}
\item ตัวเก็บประจุ 3 ตัว C ต่ออนุกรมกัน ความจุรวมเป็นเท่าใด? \hfill \textbf{\small (Extra)}
\begin{enumerate}[label=)]
\item ก) C
\item ข) C/3
\item ค) 3C
\item ง) 3C/2
\end{enumerate}
\item ตัวเก็บประจุ 100 μF ต่อกับแบตเตอรี่ 50 V ถูกตัดออกแล้วต่อขนานกับตัวเก็บประจุ 100 μF ที่ไม่มีประจุ พลังงานสูญเสียเป็นเท่าใด? \hfill \textbf{\small (Extra)}
\begin{enumerate}[label=)]
\item ก) 0.125 J
\item ข) 0.0625 J
\item ค) 0.25 J
\item ง) 0.5 J
\end{enumerate}
\item ตัวเก็บประจุแผ่นขนานที่มีประจุ Q เมื่อต่อแผ่นบวกเข้าด้วยกันและแผ่นลบเข้าด้วยกัน ผลที่เกิดขึ้นคือข้อใด? \hfill \textbf{\small (Extra)}
\begin{enumerate}[label=)]
\item ก) ประจุรวม = Q และความจุคงที่
\item ข) ประจุรวม = 2Q และความจุเพิ่มขึ้น
\item ค) ประจุรวม = 0
\item ง) ประจุรวม = 2Q และความจุลดลง
\end{enumerate}
\item ตัวเก็บประจุ 10 μF มีประจุ 10⁻³ C ถูกต่อขนานกับตัวเก็บประจุ 20 μF ที่ไม่มีประจุ พลังงานหลังต่อเป็นเท่าใด? \hfill \textbf{\small (Extra)}
\begin{enumerate}[label=)]
\item ก) 2.5×10⁻³ J
\item ข) 5×10⁻³ J
\item ค) 1×10⁻² J
\item ง) 2×10⁻² J
\end{enumerate}
\item ตัวเก็บประจุแผ่นขนานมีความจุ 8 μF เมื่อเพิ่มระยะห่างเป็น 2 เท่า ความจุจะเป็นเท่าใด? \hfill \textbf{\small (Extra)}
\begin{enumerate}[label=)]
\item ก) 2 μF
\item ข) 4 μF
\item ค) 8 μF
\item ง) 16 μF
\end{enumerate}
\item ตัวเก็บประจุ 1 μF และ 2 μF ต่อขนานกัน แล้วต่ออนุกรมกับตัวเก็บประจุ 3 μF ความจุเทียบเท่าเป็นเท่าใด? \hfill \textbf{\small (Extra)}
\begin{enumerate}[label=)]
\item ก) 1 μF
\item ข) 1.5 μF
\item ค) 2 μF
\item ง) 6 μF
\end{enumerate}
\item ตัวเก็บประจุ 4 μF ต่อกับแบตเตอรี่ 60 V แล้วตัดออก ถูกต่อขนานกับตัวเก็บประจุ 12 μF ที่มีประจุ 3×10⁻⁴ C ประจุรวมหลังต่อเป็นเท่าใด? \hfill \textbf{\small (Extra)}
\begin{enumerate}[label=)]
\item ก) 5.4×10⁻⁴ C
\item ข) 4.2×10⁻⁴ C
\item ค) 3×10⁻⁴ C
\item ง) 2.4×10⁻⁴ C
\end{enumerate}
\item ตัวเก็บประจุแผ่นขนานมีไดอิเล็กทริก εᵣ = 5 ถ้าความจุเป็น 100 μF ความจุเมื่อไม่มีไดอิเล็กทริกเป็นเท่าใด? \hfill \textbf{\small (Extra)}
\begin{enumerate}[label=)]
\item ก) 5 μF
\item ข) 20 μF
\item ค) 100 μF
\item ง) 500 μF
\end{enumerate}
\item ตัวเก็บประจุ 4 μF ต่ออนุกรมกับตัวเก็บประจุ 12 μF ต่อกับแบตเตอรี่ 48 V ความต่างศักย์คร่อมตัวเก็บประจุ 12 μF เป็นเท่าใด? \hfill \textbf{\small (Extra)}
\begin{enumerate}[label=)]
\item ก) 8 V
\item ข) 12 V
\item ค) 24 V
\item ง) 36 V
\end{enumerate}
\item ตัวเก็บประจุ 2 ตัว C ต่ออนุกรมกัน แล้วต่อกับแบตเตอรี่ V พลังงานรวมเป็นเท่าใดเมื่อเทียบกับต่อขนาน? \hfill \textbf{\small (Extra)}
\begin{enumerate}[label=)]
\item ก) 1/4 เท่า
\item ข) 1/2 เท่า
\item ค) 2 เท่า
\item ง) 4 เท่า
\end{enumerate}
\item ข้อใดไม่ใช่ปัจจัยที่มีผลต่อความจุของตัวเก็บประจุแผ่นขนาน? \hfill \textbf{\small (Extra)}
\begin{enumerate}[label=)]
\item ก) พื้นที่แผ่น
\item ข) ระยะห่างระหว่างแผ่น
\item ค) ความหนาของแผ่น
\item ง) ชนิดของไดอิเล็กทริก
\end{enumerate}
\item ตัวเก็บประจุ 2 μF ต่อกับแบตเตอรี่ 200 V มีพลังงาน U₀ ถ้าตัดแบตเตอรี่ออกแล้วเพิ่มระยะห่างเป็น 2 เท่า พลังงานจะเปลี่ยนเป็นเท่าใด? \hfill \textbf{\small (Extra)}
\begin{enumerate}[label=)]
\item ก) U₀/2
\item ข) U₀
\item ค) 2U₀
\item ง) 4U₀
\end{enumerate}
\end{enumerate}
\end{document}
